\section*{Hoofdstuk \ref{ch:b3:dna-repair} --- Genoomstabiliteit: DNA-schade, herstel en recombinatie}

\begin{solution}{exo:b3:dna-repair:1}
Uracil: spontaan (\omterm{def:b3:dna-repair:lesions}{deaminering} van C), \omterm{def:b3:dna-repair:excision}{base-excisieherstel} door het
\omterm{def:b3:dna-repair:excision}{glycosylase} voor uracil-DNA. Thyminedimeer: uit het milieu (UV),
\omterm{def:b3:dna-repair:excision}{nucleotide-excisieherstel}. \omterm{def:b3:dna-repair:lesions}{8-oxoguanine}: spontaan (oxidatie),
\omterm{def:b3:dna-repair:excision}{base-excisieherstel} door het 8-oxoG-glycosylase. G--T-mismatch achter
de vork: replicatiefout, \omterm{prop:b3:dna-repair:fidelity}{mismatchherstel}. \omterm{def:b3:dna-repair:lesions}{Dubbelstrengsbreuk} in G1:
uit het milieu (straling) of spontaan; \omterm{def:b3:dna-repair:dsb}{niet-homologe eindverbinding},
omdat er geen zusterchromatide beschikbaar is.
\end{solution}

\begin{solution}{exo:b3:dna-repair:2}
De nucleotide-excisie herkent een fysische eigenschap die alle
volumineuze \omterm{def:b3:dna-repair:lesions}{laesies} delen --- de vervorming van de helix, of een
gestrand RNA-polymerase --- en haalt een plek weg met daarin wat haar
ook veroorzaakte. Kleine \omterm{def:b3:dna-repair:lesions}{laesies} (uracil, 8-oxoG, gealkyleerde basen)
vervormen de helix nauwelijks, dus moeten zij aan hun chemie worden
herkend, met één \omterm{def:b3:dna-repair:excision}{glycosylase} per soort verkeerde base.
\end{solution}

\begin{solution}{exo:b3:dna-repair:3}
Een mismatch bestaat uit twee normale basen, en alleen die op de
nieuwe streng is fout; het systeem moet weten welke streng nieuw is.
\emph{E.~coli} merkt de ouderlijke streng met methylgroepen op GATC,
die de nieuwe streng enkele minuten lang mist; MutH knipt de
ongemethyleerde streng in. In een \emph{dam}-mutant is geen van beide
strengen gemethyleerd: MutH knipt willekeurig een van beide in, het
herstel haalt de juiste base even vaak weg als de verkeerde --- zodat
de helft van de fouten als mutatie wordt vastgelegd (een mutator) ---
en inknipjes op beide strengen bij nabije GATC-plaatsen maken
\omterm{def:b3:dna-repair:lesions}{dubbelstrengsbreuken}.
\end{solution}

\begin{solution}{exo:b3:dna-repair:4}
Ongeveer $35\times 2 = 70$ \omterm{def:b3:dna-repair:lesions}{dubbelstrengsbreuken}. Na een uur is het
merendeel samengevoegd: met een halveringstijd van bijna \qty{30}{min}
ongeveer $70/4 \approx 18$ foci. Na zes uur is vrijwel alles hersteld
($70/2^{12} < 1$) en blijven alleen enkele hardnekkige, complexe
breuken over.
\end{solution}

\begin{solution}{exo:b3:dna-repair:5}
(a) $10^{-5}\times 10^{-2}\times 10^{-3} = 10^{-10}$; $6.4\times
10^{9}\times 10^{-10} = 0.64$ mutaties per deling. (b) Zonder
\omterm{prop:b3:dna-repair:fidelity}{mismatchherstel} $10^{-7}$; $640$ per deling. (c) Zonder proeflezen
$10^{-8}$; $64$ per deling. De mutant zonder \omterm{prop:b3:dna-repair:fidelity}{mismatchherstel} is de
gevaarlijkste, duizendvoudig boven normaal.
\end{solution}

\begin{solution}{exo:b3:dna-repair:6}
$k = \ln 2/\qty{0.25}{h} = \qty{2.8}{h^{-1}}$; $D = 10^{4}/24 =
\qty{420}{h^{-1}}$; $L^{*} = 420/2.8 \approx 150$ abasische plaatsen op
enig moment aanwezig. Tienvoudig trager herstel: $L^{*} \approx \num{1500}$.
De vork passeert elke locus eenmaal en treft daar de \omterm{def:b3:dna-repair:lesions}{laesies} aan die er
op dat moment zijn; over het genoom gemiddeld komt zij er ongeveer
$L^{*}$ tegen --- zo'n $150$ normaal, $\num{1500}$ met het middel.
\end{solution}

\begin{solution}{exo:b3:dna-repair:7}
$k_{\text{NHEJ}} = \ln 2/0.5 = \qty{1.39}{h^{-1}}$, $k_{\text{HR}} =
\ln 2/4 = \qty{0.17}{h^{-1}}$; deel via HR $0.17/1.56 = 0.11$;
halveringstijd $\ln 2/1.56 = \qty{27}{min}$. Een ingestorte vork geeft
een breuk met \emph{één uiteinde}: er is geen tweede uiteinde dat NHEJ
kan samenvoegen, dus kan alleen recombinatie met de zuster de vork
herstellen --- en zij is hoe dan ook de enige nauwkeurige route.
\end{solution}

\begin{solution}{exo:b3:dna-repair:8}
(a) Lynch: instabiele microsatellieten (de lengtes verschillen tussen
tumorcellen), kanker van dikke darm, baarmoederslijmvlies, maag en
eierstok, normale reactie op zonlicht. (b) XP-A: stabiele
microsatellieten, extreme zonnebrand bij minimale blootstelling, huid-
en oogkanker in de kindertijd, en in deze groep voortschrijdende
neurologische aftakeling. (c) XP-V: stabiele microsatellieten, normaal
excisieherstel zodat de zonnebrand milder is, maar de dimeren die aan
de excisie ontsnappen worden door foutgevoelige polymerasen overbrugd
--- huidkankers, die later beginnen dan in groep~A.
\end{solution}

\begin{solution}{exo:b3:dna-repair:9}
Levercellen: Cre heeft exon~3 uit beide allelen gesneden (als een
cirkel die verloren gaat), zodat het gen vanaf de geboorte in elke
hepatocyt inactief is --- een leverspecifieke knockout. Neuronen: geen
Cre, de gefloxte allelen zijn intact en werkzaam. Tegengestelde
richting: Cre draait het exon om en weer terug, zonder ophouden, zodat
op elk moment ongeveer de helft van de cellen het omgekeerde,
niet-werkzame exon draagt --- een mozaïek, geen zuivere deletie.
\end{solution}

\begin{solution}{exo:b3:dna-repair:10}
Naarmate LexA wordt geknipt daalt zijn concentratie geleidelijk.
Operatoren die LexA zwak binden komen bij een kleine daling al vrij ---
\emph{uvrA}, \emph{uvrB}, \emph{recA} --- zodat het nauwkeurige herstel
het eerst begint; de operatoren van \emph{sulA} en van de
foutgevoelige polymerasen binden LexA stevig en komen pas vrij wanneer
hij bijna weg is, dat wil zeggen wanneer de schade lang heeft
aangehouden. Een bacterie zonder de foutgevoelige polymerasen zou
sterven bij vorken die worden geblokkeerd door \omterm{def:b3:dna-repair:lesions}{laesies} die zij niet kan
overbruggen, en zou, zonder geïnduceerde mutagenese, veel trager
resistentie tegen het antibioticum ontwikkelen --- een nadeel voor de
bacterie, een voordeel voor de patiënt; remmers van de \omterm{prop:b3:dna-repair:sos}{SOS-respons}
worden om die reden onderzocht.
\end{solution}

\begin{solution}{exo:b3:dna-repair:11}
Keten: PARP geremd $\to$ enkelstrengsbreuken blijven bestaan $\to$ een
replicatievork zet elk daarvan om in een \omterm{def:b3:dna-repair:lesions}{dubbelstrengsbreuk} met één
uiteinde $\to$ het herstel vereist \omterm{def:b3:dna-repair:dsb}{homologe recombinatie} $\to$ de
tumorcel, zonder \omterm{def:b3:dna-repair:dsb}{BRCA}, kan dat niet $\to$ verkeerd samenvoegen, verlies
van chromosomen, dood. Normale cellen, met één werkzaam allel,
recombineren en overleven. Resistentie: (1) een tweede mutatie in
\emph{\omterm{def:b3:dna-repair:dsb}{BRCA}} die het leesraam en een werkend eiwit herstelt --- aan te
tonen door de tumor of het circulerende tumor-DNA in het bloed te
sequencen, en door de terugkeer van Rad51-foci na bestraling; (2)
verlies van de eiwitten die het wegsnoeien van uiteinden blokkeren
(53BP1--shieldine), wat de recombinatie zonder BRCA1 deels herstelt ---
aan te tonen door hun afwezigheid bij kleuring en door herstelde
Rad51-foci; verder ook mutaties van PARP1 die verhinderen dat het op
het DNA blijft haken, of pompen die het middel uitscheiden, aan te
tonen door respectievelijk sequencing en expressie.
\end{solution}

\begin{solution}{exo:b3:dna-repair:12}
Twee homologe duplexen, de ene met $A$ en de andere met $a$, wisselen
over de merker heen enkele strengen uit: elk heeft nu over dat traject
één streng van elke ouder --- een heteroduplex met een
$A/a$-mismatch. Zet het \omterm{prop:b3:dna-repair:fidelity}{mismatchherstel} de mismatch op de binnengedrongen
duplex om naar $A$, dan dragen drie chromatiden $A$ en één $a$:
$3{:}1$; omgezet naar $a$: $1{:}3$; blijft zij onhersteld, dan splitst
de chromatide $A$ en $a$ bij de eerste mitose na de meiose en ontstaan
twee verschillende dochtersporen (postmeiotische splitsing, $5{:}3$ in
een ascus met acht sporen). Het heteroduplextraject beslaat maar enkele
honderden tot enkele duizenden basenparen rond de breuk die het
initieerde, dus kunnen alleen merkers daarbinnen worden omgezet, en
die liggen per constructie naast de crossover.
\end{solution}

\begin{solution}{pb:b3:dna-repair:1}
\textbf{1.} $10^{4} + 400 + 10^{4} \approx 2\times 10^{4}$ \omterm{def:b3:dna-repair:lesions}{laesies} per
dag, $7\times 10^{6}$ per jaar --- ruwweg één per kilobase van het
genoom per jaar.
\textbf{2.} $10^{-5}\times 10^{-2}\times 10^{-3} = 10^{-10}$ per bp;
$6.4\times 10^{9}\times 10^{-10} = 0.64$ mutaties per deling.
\textbf{3.} $0.64\times 0.015 \approx 0.01$ coderende mutaties per
deling; over $50$ delingen ongeveer $0.5$.
\textbf{4.} $10^{-7}$ per bp, $640$ mutaties per deling. Uitgegleden
cellen: zonder herstel $10^{-4}\times 40\times 10^{9} = 4\times 10^{6}$;
met herstel $10^{-7}\times 40\times 10^{9} = 4000$ --- een duizendvoudig
verschil, zichtbaar als \omterm{def:b3:dna-repair:mmr}{microsatellietinstabiliteit}.
\textbf{5.} Uracil hoort niet in DNA thuis, wordt herkend door het
uracil-DNA-glycosylase en uitgesneden: hersteld. Thymine uit
5-methylcytosine is een normale base tegenover een G; buiten de
replicatie markeert niets haar als fout, en zij wordt een
C$\to$T-mutatie --- het mechanisme dat \omterm{def:b3:chromatin-epigenetics:methylation}{CpG} uit het genoom heeft
uitgedund.
\textbf{6.} Het 8-oxoG-glycosylase haalt de geoxideerde base weg zolang
zij nog tegenover C staat (vóór de replicatie); het tweede \omterm{def:b3:dna-repair:excision}{glycosylase}
haalt een A weg die verkeerd tegenover 8-oxoG is ingebouwd (na de
replicatie, wat een nieuwe kans geeft om de G te herstellen); het
hydrolase vernietigt geoxideerd dGTP zodat het niet tegenover A wordt
ingebouwd. Falen alle drie, dan paart 8-oxoG met A en legt de volgende
ronde een transversie G$\cdot$C$\to$T$\cdot$A vast.
\textbf{7.} Normaal $k = \ln 2/2 = \qty{0.35}{h^{-1}}$; XP $k = \ln 2/70
= \qty{0.0099}{h^{-1}}$.
\textbf{8.} Normaal $10^{5} e^{-2.77} \approx 6200$; XP $10^{5}
e^{-0.079} \approx \num{92000}$.
\textbf{9.} Mutaties $6.2$ en $92$ per cel; coderend $0.09$ en $1.4$.
\textbf{10.} $500$ blootstellingen: $47$ coderende mutaties (normaal) en
$690$ (XP), tegen $0.5$ uit de replicatie --- zonlicht overheerst de
mutatielast van de huid zelfs bij een normaal persoon.
\textbf{11.} Coderende sequentie $9.6\times 10^{7}$ bp; de
aandrijvende genen zijn $4\times 10^{4}/9.6\times 10^{7} = 4.2\times 10^{-4}$
daarvan. Gemiddeld aantal treffers $690\times 4.2\times 10^{-4} = 0.29$;
$P(\ge 1) = 1 - e^{-0.29} =
0.25$. Van $10^{9}$ cellen dragen er $2.5\times 10^{8}$ een
aandrijvende mutatie (normaal kind: gemiddeld $0.02$, \qty{2}{\%}).
\textbf{12.} Een dimeer is geen mutatie; hij wordt er pas een wanneer
een vork hem met foutgevoelige overbrugging kopieert, dus zijn delende
cellen met veel onherstelde dimeren degene die muteren. Ultraviolet
licht bereikt alleen de huid en het oog; inwendige weefsels krijgen
geen dimeren en zijn bij XP vrijwel normaal.
\textbf{13.} $70$ breuken; na \qty{1}{h} met alleen NHEJ $70\times
2^{-2} \approx 18$ foci.
\textbf{14.} $k_{\text{NHEJ}} = \qty{1.39}{h^{-1}}$, $k_{\text{HR}} =
\qty{0.17}{h^{-1}}$; deel via HR $0.17/1.56 = 0.11$.
\textbf{15.} Halveringstijd $\ln 2/1.56 = \qty{27}{min}$; na \qty{1}{h}
$70\,e^{-1.56} \approx 15$ breuken.
\textbf{16.} $70\times 0.015\times 0.7 \approx 0.7$ verstoorde genen per
cel.
\textbf{17.} $n = 70$: $2415$ paren $\times 3\times 10^{-4} = 0.72$
translocaties. $n = 7$: $21\times 3\times 10^{-4} = 0.006$. Het aantal
breuken is $\propto D$, maar het aantal \emph{paren} gelijktijdige
breuken is $\propto n^{2} \propto D^{2}$.
\textbf{18.} $D = \qty{7}{h^{-1}}$, $L^{*} = 7/1.39 \approx 5$ breuken
tegelijk aanwezig: $10$ paren, $0.003$ translocaties --- tweehonderd
keer minder dan dezelfde dosis in één flits. Fractioneren laat normaal
weefsel tussen de fracties herstellen en houdt het aantal gelijktijdige
breuken klein.
\textbf{19.} Intact herstel: $e^{-2/1.5} = 0.26$ bij \qty{2}{Gy},
$e^{-4} = 0.018$ bij \qty{6}{Gy}. Zonder NHEJ: $e^{-4} = 0.018$ en
$e^{-12} = 6\times 10^{-6}$. $D_{0}$ is de dosis die gemiddeld één
dodelijke onherstelde \omterm{def:b3:dna-repair:lesions}{laesie} per cel achterlaat (overleving $1/e$).
\textbf{20.} De \qty{11}{\%} van de tweezijdige breuken in G2 die HR
nauwkeurig zou hebben hersteld gaat nu via NHEJ, en elke vorkbreuk met
één uiteinde --- die NHEJ in het geheel niet correct kan herstellen ---
wordt verkeerd samengevoegd: over de jaren stapelt het genoom deleties,
translocaties en kopieaantalveranderingen op, de kenmerkende littekens
van \emph{\omterm{def:b3:dna-repair:dsb}{BRCA}}-tumoren.
\textbf{21.} $10^{4}\times 0.01 = 100$ breuken met één uiteinde per dag.
Een breuk met één uiteinde heeft geen partner; NHEJ kan haar alleen
laten liggen of aan een willekeurig ander uiteinde koppelen, wat een
translocatie of een verlies oplevert.
\textbf{22.} Eén werkzaam \emph{BRCA2}-allel maakt genoeg eiwit voor de
recombinatie: de normale cellen herstellen hun vorkbreuken en overleven
de remmer.
\textbf{23.} Een tweede leesraamverschuiving stroomafwaarts van de
eerste herstelt het leesraam en levert een verkort maar deels werkzaam
BRCA2 op, zodat de recombinatie terugkeert en het middel niet meer
doodt. Aan te tonen door \emph{BRCA2} in de resistente tumor of in
circulerend tumor-DNA te sequencen, en functioneel door het opnieuw
verschijnen van Rad51-foci na bestraling.
\textbf{24.} Het doden vergt dat het \omterm{prop:b3:dna-repair:fidelity}{mismatchherstel} de gealkyleerde
basen aanpakt; een Lynch-tumor die het mist is tolerant --- het middel
doodt haar niet (zulke tumoren worden in plaats daarvan anders
behandeld, onder meer met immuuntherapie, waarvoor hun zware mutatielast
hen gevoelig maakt).
\textbf{25.} Ongeveer \num{92000} dimeren bij de vork in de XP-cel;
$10^{-10}$ mutaties per basenpaar per deling; HR herstelt \qty{11}{\%}
van de breuken in G2.
\end{solution}
